\documentclass{article}

\newif\ifshadedsectiontitles
\shadedsectiontitlestrue

\input{\doccommon}
\usepackage{fancyhdr}
\usepackage{lastpage}
\pagestyle{fancy}

\begin{document}

\thispagestyle{empty}
\cfoot{\thepage\ of \pageref{LastPage}}
\lhead{Data Structures and Objects}
\chead{Lab 6}
\rhead{Generic Containers}

\noindent{\Huge Data Structures and Objects\\[3pt]CSIS 3700}\\[3pt]
\emph{\large Lab 6 --- Generic Containers}\\[3pt]

\section{Goal}
Experiment with various methods for creating containers that can hold multiple
types of data within the same program.

\section{Method 1 --- Unions}
So far, you have seen a {\bf class} --- a collection of data and the code that
manipulates the data. You may have also seen a {\bf struct}, which is just the
collection of data without the code.\parend

A {\bf union} is similar to a {\bf struct}. The difference is, while a
{\bf struct} can hold many values inside of itself simultaneously, a {\bf union}
can only hold one value at a time. Consider the following {\bf struct}:
\begin{lstlisting}
struct {
  int foo;
  double bar;
};
\end{lstlisting}
\vskip6pt
This {\bf struct} can hold an integer {\bf foo} value and a double-precision
value {\bf bar} simultaneously.\parend
Now, consider the following {\bf union}:
\begin{lstlisting}
union {
  int foo;
  double bar;
};
\end{lstlisting}
\vskip6pt
This can hold either {\bf foo} or {\bf bar}, but not both simultaneously. The
{\bf struct} has enough space to hold both --- 12 bytes on most systems ---
while the {\bf union} only has enough space for the largest of its items ---
the {\bf double}, which is 8 bytes.\parend

To try this method out, you will need to modify the following files:
\begin{itemize}
  \item {\tt simpleStack.h}
  \item {\tt simpleStack.cc}
  \item {\tt stackTest1.cc}
\end{itemize}

First, modify {\tt simpleStack.h}, replacing the {\bf typedef} line with the
following:
\begin{lstlisting}
union UItem {
  int iVal;
  double dVal;
};

typedef union UItem StackType;
\end{lstlisting}
\vskip6pt
This defines the {\bf union} and creates an alias {\bf StackType}; the compiler
will replace all occurrences of {\bf StackType} with {\bf union UItem}.\parend

In {\tt stackTest1.cc}, create a single {\bf StackType} variable; add
\begin{lstlisting}
StackType digit;
\end{lstlisting}
as a local variable inside {\bf main(\,)}. Replace the {\bf push(\,)} call with
\begin{lstlisting}
digit.iVal = d;
myStack.push(digit);
\end{lstlisting}
\vskip6pt
Finally, replace the {\bf pop(\,)} and {\bf cout} lines with
\begin{lstlisting}
digit = myStack.pop();
cout << digit.iVal << endl;
\end{lstlisting}
\vskip6pt
Compile {\tt stackTest1.cc} and {\tt simpleStack.cc} with one {\tt g++} command;
no {\tt Makefile} is necessary here. Test the program; it should read a number
and output the digits one at a time in the proper order.

\section{Method 2 --- Zen containers}
``Argue for your limitations, and sure enough, they're yours.'' ---
\emph{Messiah's Handbook}, from Richard Bach's \emph{Illusions}\parend

Another method for containing multiple data types is to create a container that
stores no data type. Some containers, such as stacks, queues and linear lists,
do not place any special restrictions on the type of data stored within them. In
fact, these containers completely ignore the data type, other than to use it to
determine how many bytes to move into or out of the container. If that is the
case, then it is possible to devise a container where no type is specified; only
the number of bytes to move is given.\parend

I call this type of approach a ``zen'' data structure, as it has a distinctly
zen flavor to it: we can store any kind of data, since we specify no kind of
data.\parend

Note that some containers require some kind of data comparability; for example,
dictionaries require at least the ability to compare two keys to see if they are
the same, and binary search trees require the ability to compare two keys for
equality and less or greater than. Zen data structures can be adapted to work
with these (\emph{cf} the {\bf qsort(\,)} function) at the expense of
complicating the code.\parend

From the user's perspective, zen data structures differ from other structures in
two main ways:
\begin{itemize}
  \item When initializing the container, the number of bytes to move must be
  given.
  \item Adding to and removing from the container require pointers to data be
  passed.
\end{itemize}

To try this out, you will work with the following files:
\begin{itemize}
  \item {\tt zenStack.h}
  \item {\tt zenStack.cc}
  \item {\tt stackTest2.cc} --- same as stackTest1.cc
\end{itemize}
\newpage
In {\tt stackTest2.cc}, change the stack declaration to
\begin{lstlisting}
ZenStack myStack(sizeof(int));
\end{lstlisting}
This determines the number of bytes in an integer, and then tells the zen stack
to move that many bytes with each {\bf push(\,)} and {\bf pop(\,)}.\parend

Next, change the {\bf push(\,)} call to
\begin{lstlisting}
myStack.push(&d);
\end{lstlisting}
\vskip6pt
Finally, change the {\bf pop(\,)} line to
\begin{lstlisting}
myStack.pop(&d);
\end{lstlisting}
\vskip6pt
Compile {\tt zenStack.cc} and {\tt stackTest2.cc} with one {\tt g++} command.
The resulting executable should work exactly as the first one did.

\section{Method 3 --- Multiple classes}
A third method is to create a separate class for each type of data the container
needs to store. Although we could actually write out multiple classes, there is
a simpler way to accomplish this --- the {\bf template}.\parend

With a template, a placeholder is used in the class definition wherever the data
type is needed. When a template container object is declared, the actual data
type is specified. The compiler then makes a behind-the-scenes copy of the
template class, with the placeholder replaced with the actual data type.\parend

One drawback of this approach is that all code must be placed inside the class.
Non-template containers are usually split into an interface file holding the
class definition and animplementation file holding the code for the class
methods. Since the compiler needs all of the class information in order to make
copies, the code must be part of the template.\parend

To try this method out, you will work with the following files:
\begin{itemize}
  \item {\tt stack.h}, which is {\tt simpleStack.h} and {\tt simpleStack.cc}
  combined into one file
  \item {\tt stackTest3.cc}, which is the same as the other two test files
\end{itemize}
In {\tt stack.h}, remove the {\bf typedef} line, and add the following line
immediately above the class line:
\begin{lstlisting}
template <class StackType>
\end{lstlisting}
This establishes the class as a template, using {\bf StackType} as a
placeholder.\parend

Next, in {\tt stackTest3.cc}, change the stack declaration to
\begin{lstlisting}
Stack<int> myStack;
\end{lstlisting}
That's all$\ldots$ compile {\tt stackTest3.cc} and it should work just as the
other two files.\parend

\section{Method 4 --- Shallow copy}
There is a fourth method to create generic containers, but we will not attempt
it in this lab.\parend

The idea is for the container to only store pointers to data elements, not the
elements themselves. This is called \emph{shallow copying}. The pointers or
references can be made generic enough to store any kind of data in the
container; one container can even store multiple types of data simultaneously
--- a limitation on template containers. However, care must be taken to ensure
that each item stored in this manner is actually a unique element and not just
multiple pointers to the same object.

\section{What to turn in}
Please turn in your modified source code.

\end{document}
